\documentclass[12pt]{article}

\usepackage{amsmath}
\usepackage{amssymb}
\usepackage[margin=1.25in]{geometry}
\usepackage{setspace}
\usepackage{natbib}
\usepackage{hyperref}
\usepackage{parskip}

\onehalfspacing

\title{Reconsidering the Lines:\\
\large Famous Reversals in Fiscal and Monetary Doctrine}

\author{Thomas J.~Sargent}
\date{March 2026}

\begin{document}

\maketitle

\begin{abstract}
A sequel to \citet{Sargent2011} on the trade-off between price-level
stability and Ramsey-optimal tax smoothing, this essay illustrates the
practical delicacy of the institutional choice between fiscal dominance
and monetary dominance by cataloguing celebrated reversals.  Milton
Friedman crossed the divide between his 1948 proposal for coordinating
monetary and fiscal policy and his 1960 program for a fixed money-growth
rule.  Henry Clay and James Madison each opposed the First Bank of the
United States on constitutional and republican grounds, then became
architects of the Second Bank once the War of 1812 had exposed the
fiscal cost of having no national monetary institution.  Most
dramatically, Salmon P.\ Chase spent a career as a hard-money Jacksonian
Democrat, engineered the Legal Tender Acts as Secretary of the Treasury
in the Civil War, and then, as Chief Justice, wrote the opinion in
\textit{Hepburn v.\ Griswold} (1870) declaring his own Acts
unconstitutional.  The essay then turns from individual reversals to an
institutional one at the congressional level: the half-century redrawing
of the line between \emph{money} and \emph{credit}, from the bond-backed
currency of the National Banking Acts that Chase championed to the real
bills provisions of the Federal Reserve Act of 1913---a reversal that
resolved one set of problems while embedding a different instability into
the new system.  These reversals are not signs of weakness or
inconsistency; they are evidence of how contextually sensitive, and how
consequential, the institutional lines really are.
\end{abstract}

\tableofcontents
\bigskip

%=======================================================================
\section{Introduction}
%=======================================================================

Where to draw the institutional boundary between monetary and fiscal
authority is one of the deepest questions in macroeconomic design.
\citet{Sargent2011} argued that this question poses a genuine dilemma:
arrangements that support price-level stability tend to foreclose the
Ramsey-optimal tax smoothing that \citet{LucasStokey1983} showed
requires using money creation as a flexible fiscal instrument, while
arrangements that enable efficient public finance tend to inject fiscal
pressures into monetary policy and thereby risk inflation.  There is no
dominant option.  The choice depends on the fiscal environment,
the quality of monetary institutions, the credibility of commitments,
and the weight placed on different social objectives.

If the choice were straightforward, one would expect thoughtful and
principled people to make it once, correctly, and stick to it.  The
historical record suggests otherwise.  Some of the most able statesmen
and economists in American history changed their minds, sometimes more
than once, and sometimes in ways that placed them in direct tension with
their earlier selves.  This essay examines four such individual
reversals---Friedman, Clay, Madison, Chase---and then turns to a fifth
reversal that unfolded at the institutional level over half a century:
Congress's stepwise redrawing of the line between money and credit,
from the bond-backed National Bank Notes of the 1860s to the
real-bills-backed Federal Reserve notes of 1913.

Milton Friedman, who serves as the primary example in \citet{Sargent2011},
moved from a 1948 proposal that coordinated money creation tightly with
fiscal deficits---an arrangement that prioritizes efficient stabilization
but makes monetary policy subordinate to budget outcomes---to the 1960
program for a fixed $k$-percent money-growth rule that insulates money
entirely from the budget, forfeiting tax-smoothing efficiency in the
name of price-level stability.  Intellectual development accounts for part
of Friedman's shift; but the shift also reflects a judgment that the
institutional capacity for responsible fiscal coordination had been
permanently compromised by Keynesian intellectual fashions and political
incentives.

Henry Clay and James Madison present a complementary pattern at an
earlier moment in American history.  Both opposed the First Bank of the
United States---Clay as a young Jeffersonian congressman on republican
and states'-rights grounds, Madison as the leading constitutional
theorist of the founding generation.  The financial catastrophe of the
War of 1812, fought without a national bank capable of mobilizing
government credit, chastened them both.  Clay would make a national bank
the centerpiece of his``American System''; Madison would sign the bill
creating the Second Bank in 1816 having argued with new sophistication
that long legislative practice had itself resolved the constitutional
question.

The most dramatic reversal belongs to Salmon P.\ Chase.  A
hard-money Jacksonian Democrat whose early career was defined by
opposition to paper currency and to the mixing of public and private
credit, Chase became Abraham Lincoln's Secretary of the Treasury in 1861
and proceeded to design and defend the Legal Tender Acts of 1862 and
1863, which made inconvertible paper notes (\$450 million in
``greenbacks'') legal tender for all debts.  Appointed Chief Justice
in 1864, Chase spent the rest of his career trying to undo his own
handiwork.  In \textit{Hepburn v.\ Griswold} (1870) he wrote the
majority opinion holding the Legal Tender Acts unconstitutional as
applied to debts contracted before their passage.  Grant's addition of
two new justices reversed that decision in \textit{Knox v.\ Lee} (1871).
Chase dissented.

The argument is not that these figures were confused or hypocritical.
Each reversal was accompanied by sustained reasoning and, usually, by
an explicit account of why changed circumstances or changed understanding
warranted a different institutional line.  What the reversals collectively
demonstrate is that the stability-versus-efficiency trade-off identified
in \citet{Sargent2011} is not merely a theoretical curiosity; it maps
onto live political and constitutional choices that intelligent people
answer differently in different contexts, and sometimes answer differently
over the course of a single lifetime.


%=======================================================================
\section{The Framework Recalled}
%=======================================================================

\citet{Sargent2011} builds on the \citet{SargentWallace1981}
``unpleasant monetarist arithmetic'' framework.  Let the government's
flow budget constraint be
\begin{equation}
\label{eq:flow}
b_t + \tau_t + \mu_t m_{t-1} = g_t + (1+r_{t-1})\,b_{t-1},
\end{equation}
where $b_t$ is real debt, $\tau_t$ real tax revenues, $g_t$ real spending,
$r_{t-1}$ the real interest rate on debt, and $\mu_t = (M_t/M_{t-1}-1)$
the money-growth rate applied to a real money base $m_{t-1}$.  Seigniorage
$\mu_t m_{t-1}$ enters as a source of revenue that competes with ordinary
taxation.

Two polar institutional arrangements are possible:

\textit{Monetary dominance.}  The monetary authority commits credibly to
a path for $\{M_t\}$ (or equivalently the price level), and the fiscal
authority must choose $\{\tau_t,\, b_t,\, g_t\}$ to satisfy
\eqref{eq:flow} without altering that path.  Fiscal policy is the
``residual claimant''; seigniorage is fixed exogenously by the monetary
rule.  Price-level stability is guaranteed by the commitment, but
tax-smoothing efficiency is constrained: the fiscal authority cannot use
seigniorage as a flexible instrument to smooth distortionary taxes across
states and dates.

\textit{Fiscal dominance.}  The fiscal authority sets
$\{\tau_t,\, g_t,\, b_t\}$ autonomously and the monetary authority
passively accommodates whatever seigniorage is required to clear
\eqref{eq:flow}.  Ramsey-optimal tax smoothing is now achievable:
if shocks to $g_t$ are sufficiently persistent, \citet{LucasStokey1983}
show that optimality calls for using money creation alongside other taxes.
But fiscal dominance makes the price level a fiscal residual; any
deterioration in fiscal discipline is monetized and price-level
instability follows.

The core result of \citet{Sargent2011} is that neither arrangement Pareto
dominates the other.  Monetary dominance trades efficiency for stability;
fiscal dominance trades stability for efficiency.  The social ranking
depends on parameterization and context.  This is the sense in which the
``line'' is not obvious: reasonable people, reasoning carefully, can end
up on different sides depending on their assessment of the relevant
fiscal and monetary environment.  The historical reversals documented
below are vivid illustrations.


%=======================================================================
\section{Milton Friedman: From Fiscal Coordination to Monetary Rules}
%=======================================================================

Friedman's 1948 paper in the \textit{American Economic Review},
``A Monetary and Fiscal Framework for Economic Stability''
\citep{Friedman1948}, proposed an institutional architecture in which
monetary policy was explicitly subordinated to fiscal outcomes.  The
proposal rested on three pillars: (1) a requirement that commercial banks
hold 100 percent reserves against deposits, so that money creation was
entirely a government function; (2) automatic fiscal stabilizers
(progressive taxation and transfer payments) that would cause the budget
to move into surplus during booms and deficit during recessions without
discretionary action; and (3) a rule specifying that any deficit be
financed exclusively by issuing money, while any surplus be used to
retire money.  Under this system, the supply of money was mechanically
linked to the fiscal balance: the government literally printed money
to finance shortfalls and destroyed money when revenues exceeded
expenditures.

This is a case of programmatic fiscal dominance.  The monetary
instrument---the money supply---is not controlled by an independent
central bank pursuing a price-level target; it is set residually to
satisfy a fiscal accounting identity.  In the terminology of
\citet{Sargent2011}, the monetary authority has been dissolved into
the fiscal authority.  Friedman embraced this design in 1948 because he
believed that the economic cycle was predominantly driven by changes in
aggregate demand, that discretionary monetary and fiscal management were
dangerously unreliable, and that the automatic, formula-driven coupling
of money creation to fiscal outcomes would deliver macroeconomic
stabilization without giving policymakers enough discretion to do harm.
His touchstone at this stage was the \textit{coordination} of the two
policy arms, not their separation.

By the late 1950s Friedman had moved decisively.  His Fordham University
lectures, published as \textit{A Program for Monetary Stability}
\citep{Friedman1960}, advocated for a fixed $k$-percent annual growth
rate in a well-defined monetary aggregate, entirely independent of
fiscal conditions.  The Federal Reserve should expand the money supply
at a constant rate---Friedman suggested roughly 3 to 5 percent per
year---regardless of whether the budget was in surplus or deficit.
This is a textbook case of monetary dominance: the monetary rule is the
binding commitment, and the fiscal authority must accommodate it.

The intellectual distance between the two positions is large.  Under
the 1948 framework, a persistent fiscal deficit would automatically
generate money creation and hence inflation; the 1960 framework
explicitly blocks this channel.  What changed?

Friedman identified several reasons \citep{Friedman1960}, Chapter 4.
First, he had become convinced by the historical evidence, which he and
\citet{FriedmanSchwartz1963} were then assembling, that autonomous
monetary shocks---changes in the money stock not driven by fiscal
conditions---were the dominant source of macroeconomic fluctuations.
The Great Depression was not a collapse of aggregate demand correctable
by fiscal-monetary coordination; it was a catastrophic monetary
contraction engineered by Federal Reserve mistakes.  If monetary shocks
were the main danger, the right institutional response was to bind the
monetary authority, not to couple it to fiscal outcomes.

Second, Friedman had concluded that fiscal discipline in a democratic
polity was essentially unenforceable.  Governments that could finance
deficits by money creation would reliably do so.  The 1948 design assumed
that the fiscal stabilizers would be genuinely automatic and symmetric:
surpluses in good times would destroy money just as deficits in bad times
created it.  But the historical record suggested that democracies enlarged
their fiscal commitments during downturns and failed to retire them
during expansions.  A design that coupled monetary policy to fiscal
outcomes would therefore degenerate into a one-way ratchet toward
inflation.

Third, and most analytically interesting from the perspective of
\citet{Sargent2011}, Friedman had implicitly revised his assessment of
the stability-efficiency trade-off.  The 1948 framework prized efficiency
(smooth automatic stabilization) and accepted the price-level risk
entailed by fiscal coupling.  The 1960 program prized stability
(a committed monetary rule) and accepted the efficiency loss from
foreclosing flexible seigniorage.  The change was not a logical
correction of an earlier error; it was a shift in the relative weight
placed on the two objectives, driven by an updated reading of what
the American political economy could actually sustain.

Friedman's reversal is a clean illustration of the paper's main thesis:
the optimal institutional line between fiscal and monetary authority
depends on empirical and institutional features of the environment that
are themselves changing and contestable.


%=======================================================================
\section{Henry Clay: From Republican Critic to Champion of the American
         System}
%=======================================================================

Henry Clay entered Congress as a Kentucky senator in 1806, steeped in
the Jeffersonian Republican tradition that regarded the First Bank of
the United States as an unconstitutional concentration of financial power
in private hands.  When the First Bank's twenty-year charter expired in
1811 and recharter was debated, Clay---now a newly elected member of the
House---was among those who voted against it.  The recharter bill died
narrowly: the Senate divided 17--17, and Vice President George Clinton
cast the deciding vote against; the House defeated it 65 to 64.  Clay's
opposition was rooted in the Jeffersonian political economy that Andrew
Jackson would later systematize: the Bank was a creature of Hamilton's
Federalist designs, an engine of privilege and corruption, constitutionally
suspect, and dangerous to the democratic republic.

Five years later, Clay was one of the most forceful advocates in Congress
for creating the Second Bank of the United States.  The Second Bank was
chartered in April 1816 with a twenty-year term, capitalized at \$35
million (more than three times the First Bank's capitalization), and
endowed with explicit authority to issue notes that could serve as
currency nationwide.  Clay's support was enthusiastic and, more
importantly, principled: he incorporated the national bank into his
``American System,'' the integrated program of protective tariffs,
federally financed internal improvements, and a national monetary
institution that he argued was necessary for American industrial
development, national security, and economic integration.

What happened between 1811 and 1816 was the War of 1812.  In the
absence of a national bank, the Madison administration found it nearly
impossible to finance the war.  State banks proliferated and issued paper
with no common standard; the government was reduced to borrowing at
punishing terms from private financiers; military operations were
hampered by the Treasury's inability to transfer funds efficiently across
the continent.  In 1814, state banks outside New England suspended specie
payments---an event that dramatized the instability of a purely
decentralized monetary system.  The government's own fiscal operations
were threatened: Treasury notes circulated at steep discounts, and the
Secretary of the Treasury reported that he could not reliably meet
the government's payroll.

Clay's reversal is intelligible as a response to this evidence.  Before
the War of 1812, the costs of a national monetary institution---the
constitutional objections, the concentration of financial privilege, the
threat to state banking interests---loomed large relative to its
benefits, which seemed speculative and ideological.  After the war, the
corresponding cost-benefit calculation ran the other way: the
demonstrated fiscal cost of the absence of a national institution was
concrete and severe, while the constitutional objections had been blunted
by the practical success of the First Bank during its twenty years of
operation.  The line between efficient national finance and the dangers
of centralized monetary power had moved, not because the underlying
principle had changed, but because the evidence on the costs of each
institutional arrangement had accumulated.

In the framework of \citet{Sargent2011}, the First Bank had been an
instrument of fiscal-monetary coordination: it held government deposits,
facilitated tax collection, managed government debt, and provided a
measure of monetary discipline through its willingness to redeem its
notes.  The two decades after its expiration amounted to an unplanned
experiment in the absence of such coordination.  Clay read the
experimental results and updated accordingly.  That is not inconsistency;
that is rationality.

The sequel to Clay's conversion is worth noting.  When Andrew Jackson
vetoed the recharter of the Second Bank in 1832 and destroyed it by
withdrawing government deposits in 1833--34, Clay emerged as the leading
opponent of Jackson's action.  The Bank War became the defining issue
of Clay's Whig Party.  The man who had voted against the First Bank in
1811 became, by the 1830s, the foremost defender of the principle of a
national monetary institution against Jacksonian hard-money democracy.
The reversal had become the fixed point from which he did not thereafter
move.


%=======================================================================
\section{James Madison: Constitutional Principle and the Doctrine of
         Liquidation}
%=======================================================================

James Madison's trajectory runs parallel to Clay's and, in some respects,
deeper, because Madison's original opposition to the First Bank was not
merely political but philosophical---grounded in the most careful
constitutional theory of his generation.

When Alexander Hamilton proposed the First Bank in January 1791,
Madison led the House opposition.  His speech of February 2, 1791 is one
of the classic statements of strict constructionism.  The Constitution
granted Congress enumerated powers.  The power to charter a corporation
was not among them.  The ``necessary and proper'' clause could not
be interpreted to extend congressional power to any means that might be
deemed convenient or useful; ``necessary'' had to mean genuinely
indispensable to the execution of an enumerated power, and chartering a
bank was not indispensable to anything the Constitution explicitly
authorized.  Madison also invoked what would later be called
the principle of original understanding: the framers'
intentions, as manifested in the ratification debates, gave no
support for so broad an implied power.

Hamilton's counter-argument won: Washington accepted Hamilton's opinion
and signed the bill.  But Madison's position was not merely a tactical
loss; it established the intellectual framework that Jeffersonian
Republicans would invoke for the next two decades, and that Jackson's
Democrats would recycle in the Bank War of the 1830s.

In 1815, near the end of the war and his presidency, Madison vetoed a
bank bill---not on constitutional grounds, but because the proposed
bank lacked features that would make it an effective fiscal instrument
for the government.  The veto message explicitly declined to reopen the
constitutional question, which Madison said had been settled by
congressional and executive practice over twenty-four years.  The
following year, in 1816, Madison signed the bill creating the Second
Bank of the United States.

How did the author of the 1791 speech against bank constitutionality
come to sign the bill creating a larger successor?  Madison's answer
was sophisticated and, to his critics, infuriatingly convenient.
He had articulated it already in the 1815 veto message: while his own
constitutional views had not changed, the Constitution's meaning was
not fixed solely by the original understanding of individual framers.
It was also determined by the accumulated legislative, executive, and
judicial practice of the republic.  A course of action repeatedly
approved by Congress, acquiesced in by successive presidents, and
never invalidated by the courts thereby acquired constitutional
legitimacy through what Madison called the process of
``liquidation''---the settlement of constitutional meaning by
institutional practice over time.  By 1816, twenty-four years of the
First Bank's operation had, in Madison's view, liquidated the
constitutional question in favor of congressional authority.

The liquidation doctrine is philosophically consequential.  It implies
that the institutional line between permissible and impermissible
concentration of fiscal-monetary power is not fixed by original text
alone; it evolves through a common-law-like process of precedent and
practice.  Institutions that were once constitutionally contested can
become constitutionally settled, and the initial objections correspondingly
lose their force.  In the terms of \citet{Sargent2011}, the boundary
between stability-promoting and efficiency-promoting arrangements is not
only analytically ambiguous; it is constitutionally labile, subject to
reinterpretation as practice accumulates.

Madison's reversal also reflects the same lesson Clay drew from 1812:
the War had revealed the fiscal consequences of the absence of a national
monetary institution.  A republic that could not finance its wars
without submitting to ruinous private financiers was not well served by
constitutional principles, however impeccably derived.  The stability
of the polity---its fiscal survival---was itself at stake; abstract
monetary-constitutional purity had proved a luxury that the republic
could not afford.

The most important contrast to Madison is his friend and rival Thomas
Jefferson, who never accepted the liquidation doctrine and to the end of
his life insisted that the Bank was unconstitutional.  Jefferson's
consistency and Madison's flexibility represent the two sides of what is
today still a live debate: whether institutional practice can and should
revise constitutional meaning, or whether original principle must
constrain practice however inconvenient the results.  That debate is,
in the end, about who gets to move the lines---and whether the lines
can legitimately be moved at all.


%=======================================================================
\section{Salmon P.~Chase: From Hard Money to Legal Tender to
         Constitutional Revision}
%=======================================================================

The most internally dramatic of the reversals documented here belongs
to Salmon P.\ Chase.  His career traces a complete arc: from principled
hard-money doctrine, through the coercion of wartime necessity, to a
judicial crusade to undo his own executive work.  Few episodes in
American constitutional history have been dissected more searchingly
than this one.  The fullest contemporary anatomy is George Bancroft's
\textit{A Plea for the Constitution of the United States Wounded in the
House of Its Guardians} \citep{Bancroft1886}, a remarkable document
that repays study here.

\subsection{A Witness of Singular Authority}

Bancroft (1800--1891) was uniquely positioned to indict the Legal Tender
Cases.  As the foremost narrative historian of the United States---his
ten-volume \textit{History of the United States} (1834--74) was the
monumental scholarly achievement of the century---he brought to monetary
constitutionalism not merely the moralism of a hard-money partisan but
the documentary authority of a man who had spent fifty years in the
primary sources of the founding era.  He had personally known James
Madison and John Adams; had served as Secretary of the Navy under Polk
and as Minister to Prussia and Germany under Johnson and Grant; and was
84 years old when, in the wake of \textit{Juilliard v.\ Greenman}
(1884), he sat down to write his \textit{Plea}.  The book is a
comprehensive historical argument that the framers at Philadelphia had
deliberately and by large majorities stripped from the federal
government any power to make paper money legal tender for private
debts---and that a sequence of Supreme Court decisions had betrayed that
design.  His central exhibit, running through every part of the work,
is Salmon P.\ Chase.

Bancroft's constitutional argument rests on what happened in Philadelphia
on August~16, 1787.  The original draft of the Constitution had given
Congress the power ``to borrow money, and emit bills on the credit of
the United States.''  Gouverneur Morris moved to strike the words ``and
emit bills.''  The motion passed 9--2.  Madison's convention notes
recorded that the majority ``thought it advisable not to insert this
power; thinking it probable that the omission would not disable the use
of public notes so far as they would be safe and proper, and would only
cut off the pretext for a paper currency, and particularly for making
the bills a tender either for public or private debts.''  Article~I,
Section~10 then explicitly forbade the states from emitting bills of
credit or making anything but gold and silver coin legal tender.  For
Bancroft, this two-sided design---the power denied to Congress by
omission, the power denied to the states by express prohibition---was
the clearest possible evidence that the framers had closed the door to
paper legal tender at every level of government.

Chase's career is, in Bancroft's telling, the story of that door being
pried open and the constitution bleeding as a result.

\subsection{The Hard-Money Democrat}

Chase began his public career in the 1830s and 1840s embedded in the
Jacksonian political tradition.  The Jacksonian hard-money creed, which
Chase shared, held that paper money not convertible to specie was a
form of fraud on holders and debtors alike; that the issuance of bank
notes by private institutions constituted an illegitimate privilege; and
that sound money meant gold and silver coin, or notes immediately and
unconditionally redeemable in them.  Chase was a Free Soiler and then a
Republican rather than a Democrat by the 1850s, but his monetary
convictions retained their Jacksonian character: he distrusted paper,
trusted specie, and consistently argued that mixing governmental and
banking functions corrupted both.

As Governor of Ohio (1856--60) and then Senator, Chase never departed
from these positions.  When Lincoln appointed him Secretary of the
Treasury in March 1861, Chase brought this monetary philosophy with him.
His first instinct was to finance the Civil War through taxation and
borrowing at market terms, without resort to inconvertible paper.  He
raised tariffs, promoted the new income tax, and placed bond issues
with Jay Cooke and other financial houses.  In his initial
communications to Congress, Chase explicitly rejected proposals for a
legal-tender paper currency.

\subsection{The Legal Tender Acts and the 5-20 Bonds}

By December 1861 the financial situation had become critical.  Banks had
suspended specie payments; the Treasury could not meet its obligations;
military expenditures were running at rates that bond markets could not
absorb at sustainable interest rates.  Chase reversed course completely.
He endorsed, designed, and lobbied vigorously for the Legal Tender Act
of February~25, 1862, which authorized the issuance of \$150~million in
Treasury notes---``United States Notes,'' quickly styled
``greenbacks''---that were legal tender for all debts, public and
private, but not convertible to specie.  Subsequent acts in July~1862
and March~1863 authorized additional issues, bringing the total to
approximately \$450~million.

Alongside the greenbacks, Chase issued the so-called 5-20 bonds to
finance the war: six-percent coupon bonds redeemable after five and
payable within twenty years, whose authorizing statute (the Act of
February~25, 1862) specified that interest would be paid in coin but
was entirely silent on principal.  Chase bridged this statutory gap with
reputational commitment, marketing the bonds to the public---through Jay
Cooke's enormous retail campaign---with the implicit understanding that
principal, too, would eventually be repaid in gold.  The commitment was
extra-statutory: it rested on Chase's personal credibility as a
hard-money man, not on any legislative guarantee.  That would matter
enormously later.

The arguments Chase offered for the Acts were purely fiscal and
pragmatic: the government faced a choice between legal tender paper and
either national bankruptcy or negotiated independence.  He invoked the
precedent of the Revolutionary War and argued that the power to issue
legal-tender notes was inherent in sovereignty and necessary to the
government's most vital function---its survival in armed conflict.  He
did not pretend that this was consistent with his earlier views; he
argued only that emergency had foreclosed the options he would have
preferred.

Bancroft does not excuse Chase for this.  The Revolutionary precedent
was, for Bancroft, precisely the wrong one to invoke: continental
dollars were the most instructive episode in all of American monetary
history, the reductio ad absurdum of paper money, and the primary
inspiration for the Philadelphia framers' decision to close the door.
To cite the Revolutionary experience as precedent was to cite the
disease as a justification for a recurrence.

In the taxonomy of \citet{Sargent2011}, the Legal Tender Acts were a
decisive move toward fiscal dominance.  The government issued an
inconvertible currency and declared it legal tender; the ``monetary
policy'' embedded in this decision---the price level would be whatever
was required to make \$450~million in greenbacks cover real government
expenditures---was entirely subordinate to fiscal necessity.  Chase had,
under pressure of circumstances, crossed from the hard-money extreme of
monetary dominance (specie standard, no fiat paper) to the
fiscal-dominance extreme (inconvertible legal-tender paper issued at the
government's discretion).

\subsection{\textit{Hepburn v.\ Griswold} (1870): Chase's First
             Reversal}

Lincoln appointed Chase Chief Justice in December~1864.  For the
remainder of his life, Chase pursued a constitutional agenda centered on
reversing the monetary consequences of his own Treasury policy.

The occasion came in \textit{Hepburn v.\ Griswold} (75 U.S.\ 603),
decided February~7, 1870.  The case involved a debt of \$11,250
contracted in 1860, before the Legal Tender Act, and tendered for
payment in 1864 with greenbacks.  The creditor, Mrs.\ Hepburn, argued
that she was entitled to gold coin equivalent, since the contract
antedated the Act that made paper legal tender.  By a 4--3 vote (two
seats were vacant), the Court held for the creditor.  Chase wrote the
majority opinion.

The \textit{Hepburn} opinion rested on two grounds.  First, Chase argued
that compelling creditors to accept depreciated paper for pre-existing
gold debts deprived them of property without due process in violation
of the Fifth Amendment.  Second, he argued that making greenbacks legal
tender for such debts was not within the Necessary and Proper Clause:
Congress could permissibly issue paper currency to finance the war, but
the further step of compelling private creditors to accept it for
pre-existing obligations was not ``necessary'' to any express power in
the required sense.  Crucially, Chase appealed not to express textual
prohibition but to the ``spirit'' of the Constitution---a methodological
choice that would shortly be turned against him.

The opinion conspicuously did not address whether Chase had been correct
in 1862 to recommend the Acts.  But the logical entailment was
inescapable: if making greenbacks legal tender for pre-existing debts
was unconstitutional, then Chase had committed a constitutional error as
Secretary.  Justice Samuel Miller's dissent noted quietly, and
devastatingly, that the majority had struck down legislation that the
Chief Justice himself had recommended and defended as a member of the
executive branch.  Miller also criticized Chase's appeal to the
``spirit'' of the Constitution as ``abstract and intangible'' reasoning
that invited an unlimited judicial veto over congressional action.

\subsection{\textit{Knox v.\ Lee} (1871): The Court-Packing Reversal}

The reversal of \textit{Hepburn} came with stunning speed and in
circumstances whose constitutional impropriety has been debated ever
since.  On the very day \textit{Hepburn} was announced---February~7,
1870---President Grant nominated two new justices to fill existing
vacancies: William Strong of Pennsylvania and Joseph Bradley of New
Jersey.  Chase believed, and said, that Grant had deliberately timed
the nominations to reconstitute the Court and obtain a reversal.  When
the Court reached its new complement of nine, the Attorney General moved
for reargument of the legal tender question.  \textit{Knox v.\ Lee} and
\textit{Parker v.\ Davis} (79 U.S.\ 457) were decided together on
May~1, 1871, by a vote of 5--4, overruling \textit{Hepburn}.

Justice Strong's majority opinion relied on the implied-powers doctrine
of \textit{McCulloch v.\ Maryland}: if Congress has the war power and
the power to borrow money, and if a legal-tender paper currency is an
appropriate means of executing those powers, the Acts are constitutional
under the Necessary and Proper Clause.  Strong's opinion included a
remarkable passage---remarkable even by the standards of intra-tribunal
comity---chiding Chase for having ``forced'' an earlier decision when
the Court was divided and on the verge of receiving new appointees.
As Bancroft observed, this was the majority lecturing the former
majority about the indecorum of the very process by which the current
majority had come to exist.

Chase dissented, joined by Justices Clifford and Field.  His dissent
insisted that the majority had abandoned the constitutional limitations
on congressional power in favor of an open-ended theory of implied
powers that left no meaningful constraint on federal action.  Field's
dissent was the most constitutionally principled of the three:
the doctrine that whatever is not expressly forbidden may be exercised
would ``change the whole character of our government'' from one of
limited, enumerated powers to one of unlimited powers in its relations
to the people.

Bancroft regarded the Court's composition question as itself
corroborating his thesis.  On a constitutional question of the first
magnitude---whether the federal government may compel private citizens
to accept depreciated paper in discharge of gold obligations---the
outcome turned not on any change in the law or the facts but on the
personnel of the Court.  This was, in constitutional terms, not a
``hard constraint'' (in the sense of a textual prohibition too clear
to circumvent) but a ``soft constraint'' that yielded to political
pressure once the political composition of the appointing authority
changed.  The same logic that had made Chase's reputational commitment
on the 5-20 bond principal extra-statutory and therefore fragile had
made the constitutional prohibition on legal tender extra-textual and
therefore fragile.

\subsection{\textit{Juilliard v.\ Greenman} (1884): Legal Tender as
             Permanent Sovereignty}

The third and final case swept away whatever remained of the
\textit{Hepburn} principle.  In \textit{Juilliard v.\ Greenman}
(110~U.S.\ 421), decided March~3, 1884, Justice Gray's majority
opinion---joined by eight of nine justices---held that the power to
declare paper money legal tender was not limited to wartime emergencies
but was a permanent attribute of national sovereignty.  The statutory
notes at issue had originally been issued during the war, then redeemed
and reissued in peacetime under the Resumption and Refunding Act of
1878; no emergency justified them.  That did not matter.  Gray wrote:

\begin{quote}
The power to make the notes of the government a legal tender in payment
of private debts being one of the powers belonging to sovereignty in
other civilized nations, and not expressly withheld from congress by
the constitution; we are irresistibly impelled to the conclusion that
the impressing upon the treasury notes of the United States the quality
of being a legal tender in payment of private debts is an appropriate
means, conducive and plainly adapted to the execution of the undoubted
powers of congress.
\end{quote}

For Bancroft, this language was ``the language of revolution.''
The argument from what ``other civilized nations'' had done had no place
in interpreting a constitution deliberately designed to withhold powers
that every European monarchy exercised.  The argument from ``not
expressly withheld'' converted a government of enumerated powers into
one of unlimited powers: if Congress may do anything not expressly
forbidden, then the framers' painstaking enumeration of powers was a
dead letter.  Moreover, the same nine justices who joined the
\textit{Juilliard} majority had unanimously declared, just two years
earlier, that ``the government of the United States is one of delegated,
limited, and enumerated powers,'' and that ``every valid act of congress
must find in the constitution some warrant for its passage.''  Bancroft
noted acidly that the \textit{Juilliard} opinion was therefore ``in
flagrant antagonism'' to the Court's own recent, unanimous, and
repeatedly given interpretations of the constitutional structure.

The sole dissenter was Justice Field---who had been in Chase's
\textit{Hepburn} majority and in the \textit{Knox} dissent, and who
now stood entirely alone.  Bancroft's rhetorical move was to insist
that Field was not a minority of one but the transmitted judgment of
the founding generation: standing beside him were ``invincible vouchers
for the rightness of his dissent''---James Wilson, Oliver Ellsworth,
and William Paterson, all three appointed to the Court by Washington
from among the Convention's own framers.  The institutional line that
the framers had drawn at Philadelphia, that Chase had tried to restore
from the bench, and that Field had defended alone, had been entirely
erased.

\subsection{Constitutional Tragedy and Soft Constraints}

Bancroft did not treat Chase's trajectory with personal relish.  He
read it as symptomatic of a constitutional culture that had lost its
moorings rather than as individual hypocrisy.  But his framing makes
the theoretical point with unusual clarity.  Chase's reputational
commitment to gold repayment of the 5-20 bond principal and his
constitutional commitment to the invalidity of legal-tender paper were
both \textit{extra-statutory}---both existed as assurances given without
explicit legislative or textual authority.  Both were undone by the same
political logic that had originally forced Chase to accept the Legal
Tender Acts as a war measure.  When the political environment changed
enough to matter, neither commitment was self-enforcing.

For the purposes of \citet{Sargent2011}, the arc of Chase's career
encapsulates the full range of the stability--efficiency trade-off.  His
early Jacksonian position was an extreme form of monetary dominance:
specie only, no government paper, the fiscal authority must adjust
entirely to monetary reality.  The Legal Tender Acts represented a
wartime capitulation to fiscal dominance: the monetary instrument
becomes whatever fiscal survival requires.  His judicial career was an
attempt to restore monetary dominance by constitutional command---to fix
in law the institutional line against which executive necessity, however
acute, could not encroach.

The \textit{Hepburn}--\textit{Knox}--\textit{Juilliard} trilogy
demonstrates that the attempt failed, and \textit{why} it failed.
The constitutional constraint Chase sought to enforce was a soft
constraint: it rested not on language expressly prohibiting the legal
tender power but on the ``spirit'' of the Constitution as Chase
interpreted it, and on the historical evidence that the Philadelphia
Convention had deliberately struck out the ``emit bills'' clause.
Bancroft thought this soft constraint should have been treated as a
hard one; the Court disagreed, and the political environment---
unwilling to foreclose a fiscal instrument it had used to win a
civil war---did not sustain the harder reading.  Institutional lines
that are not politically sustainable cannot be maintained by
constitutional command alone, and Chase's career is the clearest proof.


%=======================================================================
\section{The Money--Credit Line: From Bond-Backing to Real Bills,
         1863--1913}
\label{sec:money-credit}
%=======================================================================

The reversals catalogued above all concern the boundary between the
monetary and fiscal \emph{authorities}: who controls the supply of
money, and whether the fiscal authority may commandeer the monetary
instrument.  But there is a related and equally contested institutional
line---the line between \emph{money} and \emph{credit}---and the half-
century following Chase's National Bank Acts is a case study in
Congress repeatedly redrawing it.

\subsection{Bond-Backing and the Line at Government Debt}

The National Bank Act of 1864 drew the money--credit line at an
unusual location: United States government bonds.  National Bank
Notes---the dominant hand-to-hand currency of the post-Civil War
era---could be issued only against specified federal bonds deposited
with the Comptroller of the Currency.  That institutional officer,
whose primary function was verifying the bond portfolios, could
certify the soundness of every dollar in circulation.  The line
between what counted as ``money-eligible'' and what counted as mere
``credit'' was the difference between government bonds on one side and
every other asset---commercial loans, railroad securities, warehouse
receipts, farm mortgages---on the other.

In the framework of \citet{Sargent2011} this arrangement is a form of
\emph{monetary dominance}, but with a peculiar twist: the binding
constraint on the money supply was not a target set by an independent
central bank but the stock of outstanding federal debt.  When the
government ran the large primary surpluses of the 1870s and 1880s and
retired war bonds, the supply of eligible collateral contracted, and
the currency supply contracted with it regardless of any commercial
demand for more money.  Fiscal virtue, no less than fiscal profligacy,
could thus distort the monetary instrument; the two distorted it in
opposite directions, but neither direction was chosen for monetary
reasons.  The bond-backing requirement was a hard constraint against
inflationary over-issue, but it purchased that stability by tying the
money supply to the government's debt-management calendar rather than
to the needs of trade.

\subsection{The Inelasticity Debate as Four Decades of Disagreement
            about the Line}

Critics from the 1870s onward called the resulting currency
``inelastic.''  The agrarian South and West needed sharply higher
volumes of currency during planting and harvest seasons; the bond-backed
system offered no mechanism for such expansion.  Financial panics
punctuated the era with depressing regularity.  The political response
produced two broad coalitions, each proposing to move the
money--credit line in a different direction.

Greenbackers and silver inflationists wanted to dissolve the
line almost entirely.  If the government simply issued more fiat
currency---greenbacks to be increased rather than retired, or silver to
be coined at a generous ratio---there was no need for any specific asset
to ``back'' the money at all.  This was, at the limit, the same fiscal
dominance that Chase had temporarily embraced during the war: the money
supply would be whatever fiscal or political convenience required.
William Jennings Bryan's ``Cross of Gold'' speech of 1896 crystallized
this position in its most politically potent form.

Asset-currency reformers drew the line elsewhere.  Their complaint was
not too little money but the wrong kind of backing.  They argued that
national banks should be permitted to issue notes against a wider range
of productive assets---railroad and corporate bonds, agricultural
warehouse receipts, or any ``good assets'' on a bank's balance
sheet---rather than against government bonds alone.  This was a proposal
to relocate the money--credit boundary from ``government debt only'' to
``any sufficiently sound private claim.''  By broadening the eligibility
set, they hoped, the currency supply would be free to expand with
productive activity rather than with the Treasury's debt schedule.

The Aldrich--Vreeland Act of 1908, passed in the wake of the Panic
of~1907, was a temporary halfway house.  It allowed groups of national
banks to issue emergency currency backed by ``almost any securities the
banks were holding,'' but imposed a punitive tax schedule to ensure
rapid retirement once conditions normalized.  More consequentially, it
created the National Monetary Commission, chaired by Senator Aldrich,
whose multi-volume studies of European central banking provided the
intellectual scaffolding for the Federal Reserve Act of 1913.

\subsection{The Real Bills Resolution and Its Costs}

The Federal Reserve Act resolved the inelasticity debate by drawing
the money--credit line at a new location: short-term, self-liquidating
commercial paper.  Section~13 of the Act confined Federal Reserve
discount operations to notes and bills ``arising out of actual
commercial transactions'' maturing within ninety days, and explicitly
excluded paper drawn to finance positions in ``stocks, bonds, or other
investment securities.''  The so-called \emph{real bills doctrine}---
also called the commercial loan theory of banking, with roots in Adam
Smith's \textit{Wealth of Nations} and a long subsequent
history---became the operating principle of the new institution.

The new line distinguished money from credit not by asset class but
by the \emph{purpose and duration} of the underlying transaction.  A
sixty-day bill drawn against a shipment of grain was eligible because
the transaction was real, productive, and self-liquidating: when the
grain was sold, the bill would be repaid and the money supply would
contract automatically.  A call loan financing stock market speculation
was ineligible.  In the political rhetoric of the era the line ran
between Main Street and Wall Street; in the theoretical rhetoric it
ran between ``money backed by real production'' and ``money backed
by finance.''

The real bills criterion solved the inelasticity problem: the money
supply could now expand with productive activity rather than with the
stock of federal debt.  But it imported a different instability, one
that Henry Thornton had identified as early as 1802 in his
\textit{Paper Credit of Great Britain} \citep{Thornton1802}.  By
accommodating every credit demand arising from commercial transactions
valued at \emph{current prices}, the criterion set up a positive
feedback from prices to money creation to prices.  In a boom, prices
rise; the nominal value of commercial transactions rises; the demand for
discount rises; and the Fed, faithfully applying the real bills
criterion, expands the money supply---adding fuel to the expansion
rather than restraining it.  The criterion is inherently procyclical.

The most consequential application came in 1929--1933.  The Federal
Reserve Board member Adolph Miller launched a ``Direct Pressure'' policy:
any bank seeking discount window relief had to certify that it had never
made speculative loans.  Banks, unwilling to submit to the implied
interrogation, refrained from borrowing reserves even in the face of
withering deposit outflows.  Instead of discounting at the window, banks
failed---nearly 9{,}000 of them.  \citet{FriedmanSchwartz1963} held
the resulting one-third contraction of the money stock largely
responsible for the Great Depression.  Where the bond-backed currency
of the national banking era had been passively contractionary in the
face of fiscal surpluses, the real bills currency proved passively
contractionary in the face of banking panics.  Each line carried its
own characteristic failure mode.

\subsection{The Theoretical Afterlife and the Sargent--Wallace
            Reconsideration}

The real bills debate connects directly to the framework of
\citet{Sargent2011}.  In a 1982 article, \citet{SargentWallace1982}
showed that within overlapping-generations models the real bills
prescription---unfettered private intermediation against productive
assets---can be Pareto superior to the quantity-theory alternative,
which imposes restrictions on intermediation in order to control the
price level.  This is a version of the stability--efficiency trade-off
that \citet{Sargent2011} identifies in the fiscal--monetary domain:
the quantity-theory regime purchases price-level stability at the cost
of constrained intermediation; the real bills regime allows efficient
intermediation at the cost of weaker price-level control.
\citet{Laidler1984} objected that the Sargent--Wallace version had
abstracted away the key feature of the historical doctrine---the
restriction to specific collateral types---but the theorem demonstrated
that the money--credit distinction is not merely an institutional
accident or a historical curiosity.  It reflects a genuine analytical
trade-off whose terms depend on the structure of the economy and the
quality of the regulatory environment.

The congressional progression from bond-backed currency to real
bills---from the National Bank Acts of 1863--64 to the Federal Reserve
Act of 1913, through four decades of inelasticity complaint and
asset-currency agitation---thus illustrates in miniature the same
dynamic as the individual reversals catalogued in this paper.  The
actors who drafted the Federal Reserve Act were consciously reversing
the institutional design that Chase had supervised as Secretary of the
Treasury.  Their reversal solved one class of problems while creating
another; and the real bills doctrine embedded in the new Act would
itself require substantial revision in the Banking Acts of 1932--35,
a revision precipitated by the very failure mode that Thornton had
foreseen in 1802.  The half-century arc from National Bank Notes to
Federal Reserve notes is another turn of the same wheel: each new
institutional line is drawn in response to the failures of the previous
one, and each generates its own successor crisis.

%=======================================================================
\section{Common Themes}
%=======================================================================

Several themes recur across these reversals.

\textit{War as a forcing function.}  Three of the four reversals were
precipitated by war or its aftermath: the War of 1812 converted Clay
and Madison; the Civil War broke Chase.  War generates fiscal demands
that peaceful institutions are designed to defer or distribute across
time; when those demands are concentrated and immediate, they overwhelm
standing institutional arrangements.  The fiscal-dominance
arrangement---using monetary instruments to finance emergency
expenditures---becomes irresistible precisely because the alternative
is defeat.  This is not an argument against institutional constraints on
fiscal-monetary coordination; it is a reminder that any such constraints
must be evaluated in part by their behavior under the extreme cases that
history periodically generates.  \citet{SargentWallace1981} and
\citet{Sargent2011} focus on steady-state comparisons; the historical
record suggests that the transitions between regimes matter at least
as much.

\textit{The reflexivity of experience.}  Each of the reversals was
accompanied by an appeal to experience that the original position could
not have anticipated.  Clay and Madison could not have known in 1791 or
1811 what the War of 1812 would reveal about the fiscal consequences of
the absence of a national monetary institution.  Friedman could not have
known in 1948 what his and Schwartz's subsequent historical research
would reveal about the monetary origins of the Great Depression, or
what the postwar Keynesian political economy would reveal about the
systematic bias toward fiscal expansion.  Experience updated
the effective parameters of the trade-off identified in
\citet{Sargent2011}: the relative magnitude of the fiscal and monetary
risks, and the capacity of political institutions to manage them.

\textit{The constitutional versus the pragmatic argument.}  Chase's
career and Madison's liquidation doctrine both reveal a tension between
constitutional arguments about where institutional lines may be drawn
and pragmatic arguments about where they should be drawn.  The
constitutional argument tries to fix the line against the pressures of
the moment; the pragmatic argument allows the line to respond to
changing circumstances.  Madison resolved the tension by treating
constitutional meaning as itself partly determined by experience
(the liquidation doctrine).  Chase resolved it by insisting that the
Constitution fixed the line absolutely---and discovered that the
political economy would not sustain that absolutism.

\textit{The ratchet problem.}  Friedman's final position---the fixed
money-growth rule---was motivated in part by his assessment that the
coupling of monetary and fiscal policy in a democracy reliably moves
in one direction: toward more money creation, more inflation, more
fiscal expansion.  The historical record broadly supports this concern.
Neither Clay nor Madison, having supported the Second Bank, sought to
use it as an instrument of fiscal dominance; but Jackson's destruction
of the Bank was followed by the pet-bank era, the suspension of 1837,
and eventually the Legal Tender Acts.  The absence of a durable national
monetary institution contributed to a sequence of fiscal-monetary crises
that culminated in the most extreme form of fiscal dominance the republic
had seen.  The ratchet runs toward fiscal dominance when monetary
institutions are weak; only strong, credible monetary commitments can
reverse it.

\textit{The impossibility of personal commitment.}  Perhaps the deepest
lesson of Chase's career is that individuals cannot, through the force
of their own convictions, maintain institutional lines against political
and fiscal pressures.  Chase's personal monetary principles were
genuine and deeply held; he had stated them publicly for decades.
Under sufficient pressure---the fiscal demands of a war for national
survival---he abandoned them.  And having abandoned them, he could not
restore them from his judicial position, because the political forces
that had driven the abandonment had not themselves changed.
Institutional lines must be maintained by institutions, not by
individuals; and as Madison's liquidation doctrine suggests, institutions
themselves are mutable.


%=======================================================================
\section{Conclusion}
%=======================================================================

The trade-off between price-level stability and Ramsey-optimal fiscal
efficiency identified in \citet{Sargent2011} is not merely a property of
theoretical models.  It is a live institutional choice that has generated
passionate disagreement, principled reversal, and constitutional crisis
throughout American economic history.  The four individual cases examined
here---Friedman's shift from fiscal coordination to monetary rules; Clay's
conversion from Republican opponent to Whig champion of the national
bank; Madison's move from strict-constructionist critic to pragmatic
signer of the Second Bank bill; and Chase's extraordinary journey from
hard-money Democrat to greenback architect to constitutional
counter-revolutionist---all illustrate the same underlying truth: the
optimal institutional line responds to circumstances, and circumstances
change.

That truth is confirmed, at the legislative rather than the individual
level, by the half-century redrawing of the money--credit line examined
in the preceding section.  The bond-backed currency of the National
Banking System was a stable but inflexible answer to the question ``what
should money be backed by?''  The real bills doctrine of the Federal
Reserve Act was a more elastic answer that solved the inelasticity problem
while importing the procyclical instability that Thornton had identified
in 1802 and that Friedman and Schwartz held responsible for the
catastrophe of 1929--33.  The theoretical reconsideration by
\citet{SargentWallace1982}, showing that such a regime can be Pareto
superior to a quantity-theory alternative, demonstrates that this
congressional reversal, like the individual reversals, reflects a genuine
analytical trade-off rather than mere confusion.  Neither side of the line
dominates the other; the right answer depends on the fiscal and regulatory
environment.

This is not an argument for institutional nihilism.  The historical
record also supports Friedman's final conviction that price-level
stability requires institutional commitments that are robust to fiscal
pressure, and that such robustness is hard to achieve and easy to
destroy.  The Second Bank was chartered in 1816 in response to the
fiscal lessons of 1812 and destroyed by 1836 in response to political
pressures that the Bank's institutional designers had not foreseen.
The Legal Tender Acts were upheld in \textit{Knox v.\ Lee} by a Court
whose composition had been manipulated for the purpose and confirmed in
\textit{Juilliard v.\ Greenman} by an 8--1 majority that included no
defender of Chase's original position.  The trajectory is one of
progressive erosion of monetary constraints under fiscal pressure.

The cases do counsel, however, for humility about institutional design.
Principled people who have thought carefully about where to draw the lines
have repeatedly found themselves unable to maintain those lines in the
face of changed circumstances, new evidence, and the coercive logic of
fiscal emergency.  Designing institutions that are both responsive enough
to survive such pressures and durable enough to provide the stability
that is their raison d'\^etre is the enduring challenge.  The historical
reversals documented here are a measure of how difficult that challenge
is.

\bigskip

%=======================================================================
\bibliographystyle{plainnat}
\bibliography{where_to_draw_lines_sequel}
%=======================================================================

\end{document}
